\section{Background and Related Work}
\subsection{Keyword-Based Resume Screening}
Early ATS solutions used Boolean search and keyword frequency counts to match resumes against job descriptions. Such systems are simple to implement and can process large document batches quickly \cite{b1}, \cite{b3}. Their weakness is that they treat language as exact strings. A candidate who writes ``backend service development'' may be relevant to a backend engineering role, but a strict keyword system may not reward that expression unless it contains the exact expected keywords. Rigid matching also encourages applicant behavior that optimizes for the parser rather than for truthful skill representation.

\subsection{Classical NLP and Vector-Space Models}
Classical machine learning methods provide a middle layer between exact keyword search and deep semantic modeling. TF-IDF assigns greater importance to terms that are frequent in a document but not common across the corpus. When used with cosine similarity, it provides a fast numerical estimate of textual relevance \cite{b7}, \cite{b10}. In resume screening, TF-IDF is useful because it reduces the influence of generic terms and can penalize verbose, unrelated content. Rule-guided NLP can complement TF-IDF by extracting structured attributes such as skills and years of experience \cite{b16}.

\subsection{Transformer-Based Semantic Matching}
Transformer-based models such as BERT and Sentence-BERT improve semantic matching by embedding sentences into vector spaces where related meanings can be compared \cite{b6}, \cite{b12}. Recent recruitment systems use such embeddings, LLM feedback loops, and multi-agent frameworks to match candidate profiles and job descriptions more flexibly \cite{b2}, \cite{b5}, \cite{b11}, \cite{b13}. These models are valuable when candidates and recruiters use different vocabulary for the same competence. However, they demand more memory and computation than lexical approaches, and they may not naturally produce explanations that are easy for non-technical HR users to audit.

\subsection{Explainable AI in Recruitment}
Explainable AI is especially important in hiring because ranking decisions can affect career opportunities \cite{b8}, \cite{b14}. General-purpose methods such as SHAP and LIME can help analyze complex models, but their outputs may be too technical for everyday recruitment workflows \cite{b9}. A recruiter usually needs concrete statements such as ``matched two out of four required skills'' or ``missing PostgreSQL and REST'', not only feature-attribution plots. The present system therefore uses a rule-based explanation layer aligned with the internal score dictionary.

\subsection{Format Bias and OCR-Based Parsing}
Resume screening systems can unintentionally penalize candidates because of file format rather than content. A parser may fail on scanned PDFs, image resumes, or visually designed resumes with non-linear layouts \cite{b7}, \cite{b15}. OCR fallback reduces such format bias by recovering text when native extraction fails. Although OCR can introduce character-level noise, image preprocessing such as grayscale conversion, binarization, and sharpening can improve extraction quality for standard resume images.
